\section{Ricart--Agrawala mutual exclusion}\label{sec:ra}
Handout: LAB1 §4.4 \\
Code: \src{protocols/randa}

% ---------------------------------------------------------------------
\subsection{Algorithm}
To enter the critical section (CS), node $i$ stamps a request with its Lamport
timestamp $(L, i)$ and sends \texttt{Request} to all $n-1$ other nodes. It enters
when it has received a \texttt{Reply} from all of them. A node receiving a request:
\begin{itemize}[itemsep=1pt]
  \item replies at once if it is neither holding nor wanting the CS, or if it wants
        the CS but the incoming timestamp is smaller (higher priority);
  \item otherwise \emph{defers} the reply until it leaves the CS.
\end{itemize}
There is no \texttt{Release} message: the deferred replies act as the release.
Messages per CS entry: $2(n-1)$ (Lamport mutual exclusion needs $3(n-1)$).

\TODO{2--4 sentences in your own words: why the timestamp total order prevents
deadlock and guarantees mutual exclusion.}

% ---------------------------------------------------------------------
\subsection{Implementation}
\begin{itemize}[itemsep=2pt]
  \item Lock state machine (indifferent / waiting / holding):
        \texttt{ralock} in \src{protocols/randa/locks.go}
  \item Handling an incoming request (the ``defer'' decision): \texttt{ralock.question}
        in \src{protocols/randa/locks.go}
  \item Broadcasting requests and waiting for all replies: \texttt{ralock.lock}
  \item RPC handlers: \src{protocols/randa/application.go};
        Lamport clock: \src{protocols/randa/clock.go}
  \item \TODO{local lock / how concurrent Acquire on the same node is handled}
\end{itemize}

\phead{State fields reported}
\hint{LAB1 §4 asks all mutual-exclusion protocols to report \fld{role},
\fld{entries}, \fld{lastWaitMs}, \fld{messagesSent}. Make sure randa reports these.}
\begin{table}[H]\centering\small
\begin{tabularx}{\linewidth}{@{}lX@{}}
\toprule
Field & Meaning \\ \midrule
\fld{role}         & \TODO{(released / wanted / critical)} \\
\fld{entries}      & \TODO{} \\
\fld{lastWaitMs}   & \TODO{} \\
\fld{messagesSent} & \TODO{} \\
\fld{deferred}     & \TODO{} \\
\fld{clock}        & \TODO{} \\
\bottomrule
\end{tabularx}
\end{table}

\phead{Application inputs}
\begin{table}[H]\centering\small
\begin{tabularx}{\linewidth}{@{}llX@{}}
\toprule
Input & Fields & Effect \\ \midrule
\TODO{\fld{Application.Acquire}} & \fld{holdMs} & \TODO{} \\
\TODO{\fld{Application.Auto}}    & \TODO{}      & \TODO{} \\
\fld{Application.Lock} / \fld{Unlock} & --- & \TODO{} \\
\bottomrule
\end{tabularx}
\end{table}

% ---------------------------------------------------------------------
\subsection{Scenario coverage}
\begin{table}[H]\centering\small
\begin{tabularx}{\linewidth}{@{}llXl@{}}
\toprule
Report section & Handout & What it shows & Export \\ \midrule
\ref{sec:ra-normal}  & A         & two simultaneous requests & \TODO{file} \\
\ref{sec:ra-crash}   & B         & an uninvolved node crashes, everyone blocks & \TODO{file} \\
\ref{sec:ra-jitter}  & C         & random load, exponential jitter, FIFO vs.\ unordered & \TODO{file} \\
\ref{sec:ra-variant} & variant   & \TODO{} & \TODO{file} \\
\bottomrule
\end{tabularx}
\end{table}

% ---------------------------------------------------------------------
\scenario{Normal (handout A): two nodes request at the same time}{ra-normal}
\phead{Setup}
\begin{setuptable}
Nodes / latency & 3 nodes, \TODO{}\,ms, Docker \\
Input           & node-1 and node-3: \fld{Acquire(holdMs=}\fillin\fld{)} concurrently \\
Export          & \TODO{} \\
\end{setuptable}
\phead{Steps}
\begin{steps}
  \item \TODO{}
\end{steps}
\phead{Result}
\begin{table}[H]\centering\small
\begin{tabular}{@{}lrrrr@{}}
\toprule
Node & request timestamp & entered CS at & \fld{lastWaitMs} & \fld{messagesSent} \\ \midrule
node-1 & \fillin & \fillin & \fillin & \fillin \\
node-3 & \fillin & \fillin & \fillin & \fillin \\
\bottomrule
\end{tabular}
\end{table}
\screenshot{ra-normal.png}{Crossing requests; the higher-priority node replies immediately, the other defers.}
\phead{Explanation}
\TODO{3--5 sentences. Point to the deferred reply in the diagram. Count messages:
is it $2(n-1)$ per entry?}

% ---------------------------------------------------------------------
\scenario{Fault A (handout B): an uninvolved node crashes}{ra-crash}
\phead{Steps}
\begin{steps}
  \item Crash node-2 while it is not requesting the CS.
  \item node-1 calls \fld{Acquire}.
  \item \TODO{recover node-2 and show what happens}
\end{steps}
\screenshot{ra-crash.png}{node-1's request to node-2 is dropped; node-1 never enters.}
\phead{Explanation}
\TODO{3--5 sentences. Why does one crash block everyone? Compare with Lamport
mutual exclusion (same weakness).}

% ---------------------------------------------------------------------
\scenario{Fault B (handout C): random load on an exponentially jittered network}{ra-jitter}
\phead{Setup}
\begin{setuptable}
Nodes / latency & 5 nodes, exponential latency, mean \TODO{}\,ms \\
Load            & \fld{Auto} on all nodes, \TODO{rate / duration} \\
Ordering        & run 1: FIFO (default); run 2: \texttt{"ordering": "unordered"} \\
Export          & \TODO{} \\
\end{setuptable}
\phead{Check}
\fld{atMostOne} on \fld{role == critical} over the whole history
(\src{scripts/scenarios.mjs}): \TODO{pass/fail for each run}.
\phead{Result}
\begin{table}[H]\centering\small
\begin{tabular}{@{}lrrrl@{}}
\toprule
Ordering & CS entries & mean wait (ms) & messages / entry & mutual exclusion held? \\ \midrule
FIFO      & \fillin & \fillin & \fillin & \fillin \\
unordered & \fillin & \fillin & \fillin & \fillin \\
\bottomrule
\end{tabular}
\end{table}
\phead{Explanation}
\TODO{3--5 sentences. Does reordering break mutual exclusion? Why / why not for
\emph{this} implementation (note: RPC request and reply are paired).}

% ---------------------------------------------------------------------
\scenario{Design variant: \TODO{name}}{ra-variant}
\hint{Handout §4.3's variant is the message-count comparison: run the same scenario
on Lamport mutual exclusion (\nolinkurl{protocols/lamport_mutex}) and on R--A and
compare \fld{messagesSent} ($3(n-1)$ vs.\ $2(n-1)$). A stronger variant that shows
\emph{why} a rule is needed: (1) do not defer -- reply immediately even while
holding the CS -- and show two nodes in \fld{role=critical}; or (2) break ties by
something other than (timestamp, node id) and show a deadlock.}
\phead{Code change}
\TODO{}
\phead{Result}
\begin{table}[H]\centering\small
\begin{tabular}{@{}lrrl@{}}
\toprule
Version & $n$ & messages / CS entry & mutual exclusion held? \\ \midrule
Ricart--Agrawala (baseline) & \fillin & \fillin & \fillin \\
\TODO{variant or Lamport ME} & \fillin & \fillin & \fillin \\
\bottomrule
\end{tabular}
\end{table}
\phead{Explanation}
\TODO{3--5 sentences.}

% ---------------------------------------------------------------------
\subsection{Discussion questions}
\begin{qa}{Why does Ricart--Agrawala need only $2(n-1)$ messages while Lamport needs $3(n-1)$?}
\TODO{2--4 sentences.}
\end{qa}
\begin{qa}{Why does a crash of an uninvolved node block everyone?}
\TODO{2--4 sentences.}
\end{qa}
\begin{qa}{Does message reordering break mutual exclusion?}
\TODO{2--4 sentences.}
\end{qa}

\subsection{Boundaries (what is not implemented)}
\begin{itemize}[itemsep=1pt]
  \item \TODO{e.g.\ no failure detection, no persistence, fixed membership.}
\end{itemize}
